% Figure for Oscillations, Waves and Optics §A7.1 (single-slit intensity I/I0 = (sin(beta)/beta)^2, beta = pi D sin(theta)/lambda, computed against D sin(theta)/lambda; dashed: the same curve magnified ten times outside the central maximum, to show the secondary maxima of 4.7% and 1.6% near D sin(theta)/lambda = 1.43 and 2.46).
\begin{tikzpicture}[font=\small]
  \begin{axis}[nfig, clip=false, width=11cm, height=5.2cm, xmin=-3.6, xmax=4.4, ymin=0, ymax=1.18,
      xlabel={$D\sin\theta/\lambda$}, ylabel={$I/I_0$},
      xtick={-3,-2,-1,0,1,2,3}, ytick={0,0.5,1},
      xlabel style={at={(axis description cs:1.0,0)}, anchor=west},
      ylabel style={rotate=-90, at={(axis description cs:0,1)}, anchor=south}]
    \addplot[black!45, dashed, thin, domain=1:3.5, samples=300] {10*(sin(deg(pi*x))/(pi*x))^2};
    \addplot[black!45, dashed, thin, domain=-3.5:-1, samples=300] {10*(sin(deg(pi*x))/(pi*x))^2};
    \addplot[figR, domain=-3.5:3.5, samples=700] {(sin(deg(pi*x+1e-4))/(pi*x+1e-4))^2};
    \node[font=\scriptsize, figR, anchor=south west] at (axis cs:0.12,0.98) {central maximum};
    \node[font=\scriptsize, black!60, anchor=south] at (axis cs:1.43,0.49) {$\times 10$};
    \node[font=\scriptsize, black!60, anchor=west, align=left] at (axis cs:3.55,0.3) {zeros at\\ $D\sin\theta = m\lambda$};
  \end{axis}
\end{tikzpicture}
